线性时不变系统通常可以通过传递函数 $G(s)$ 描述。
本实验研究的基本传递环节包括 P、I、PT1 和 PT2 环节。
它们的动态特性可以通过阶跃响应、正弦响应、Bode 图以及零极点图进行分析。

\subsection{基本传递环节}

P 环节只对输入信号进行比例放大，其传递函数为
\[
  G_P(s) = K_P .
\]
因此，阶跃输入会立即产生阶跃输出；对正弦输入而言，输出幅值改变，但理想情况下不产生相位偏移。
P 环节没有动态储能元件，因此没有极点。

I 环节对输入信号积分，其传递函数为
\[
  G_I(s) = \frac{K_I}{s}.
\]
单位阶跃输入经过 I 环节后得到线性上升的斜坡信号。
该环节没有有限的稳态终值，因此其阶跃响应不能像 P 或 PT1 环节那样收敛到常数。
I 环节在原点具有一个极点，因此其输出会持续累积输入信号。

PT1 环节是一阶惯性环节，其典型传递函数为
\[
  G_{PT1}(s) = \frac{K}{T s + 1}.
\]
其中 $T$ 为时间常数。
在单位阶跃输入下，输出以指数形式接近稳态终值。
时间常数越大，系统响应越慢。
在 $t=T$ 时，PT1 环节的阶跃响应约达到稳态值的 $63.2\%$。

PT2 环节是二阶惯性环节，其标准形式可写为
\[
  G_{PT2}(s) =
  \frac{K \omega_n^2}{s^2 + 2D\omega_n s + \omega_n^2}.
\]
其中 $\omega_n$ 为固有频率，$D$ 为阻尼系数。
固有频率主要影响响应速度，阻尼系数决定系统是否发生振荡以及超调的大小。

实验任务中给出的 PT2 环节也常写成
\[
  G_{PT2}(s) = \frac{K}{T^2s^2 + 2DTs + 1}.
\]
两种写法等价，其中 $\omega_n = 1/T$。
因此，若传递函数分母为 $a s^2 + b s + 1$，则可由
\[
  T = \sqrt{a}, \qquad
  \omega_n = \frac{1}{\sqrt{a}}, \qquad
  D = \frac{b}{2\sqrt{a}}
\]
计算固有频率和阻尼系数。

\subsection{阶跃响应与正弦响应}

阶跃响应用于观察系统从一个工作点过渡到另一个工作点的过程。
对于稳定的 P、PT1 和阻尼 PT2 环节，输出最终趋近于稳态值 $G(0)$。
I 环节由于存在原点极点，对阶跃输入会产生斜坡响应，因此不存在有限稳态值。

对于正弦输入，线性时不变系统在暂态过程结束后仍输出同频率的正弦信号。
输出信号相对于输入信号只改变幅值并产生相位偏移：
\begin{align*}
  u(t) &= \hat{u}\sin(\omega t),\\
  y(t) &= \hat{u}\,|G(j\omega)|
          \sin\!\left(\omega t + \angle G(j\omega)\right).
\end{align*}
这一性质是用频率响应分析系统的基础。

\subsection{PT2 环节的极点}

PT2 环节的极点由分母多项式决定：
\[
  s_{1,2} = -D\omega_n \pm j\omega_n\sqrt{1-D^2}
  \qquad (0 \leq D < 1).
\]
当 $0 < D < 1$ 时，极点为共轭复数，系统表现出衰减振荡。
当 $D = 1$ 时，系统处于临界阻尼状态，在无超调条件下较快接近稳态值。
当 $D > 1$ 时，两个极点位于负实轴上，响应不再振荡，但通常更加迟缓。
当 $D = 0$ 时，极点位于虚轴上，系统不收敛到稳态值，而是持续振荡。

稳定性可直接由极点位置判断。
连续系统所有极点均位于左半平面时，系统渐近稳定。
若极点位于虚轴上，系统处于临界稳定状态；若存在右半平面极点，则系统不稳定。

\subsection{频率响应}

Bode 图用幅频特性和相频特性描述系统对不同频率正弦输入的响应。
对于 PT2 环节，低阻尼时可能出现明显的共振峰。
随着阻尼增大，共振峰减小，幅值曲线更加平滑。
在高频范围内，二阶系统的相位最终趋近于 $-180^\circ$。
幅值通常以 dB 表示：
\[
  L(\omega) = 20\log_{10}|G(j\omega)|.
\]
由于横轴采用对数频率刻度，Bode 图能够在较宽频率范围内紧凑显示系统特性。
因此，Bode 图也被称为频率响应图。
